\section{The linearised solution and the nonlinear welfare check}\label{app:nonlin}

This appendix collects two robustness exhibits that support the main-text
results without being part of the argument: the accuracy of the linearised
(sequence-space) solution given the discrete adoption choice, and the
first-order welfare check on the nonlinear transition.

\subsection{Accuracy of the linearised solution}
\label{sec:nonlin}

A natural concern is that the discrete adoption choice makes the economy too
nonlinear for the linearised (sequence-space) solution to be reliable, in
contrast to \citet{auclert2023energy}, who find nonlinearity immaterial in the continuous
economy. Comparing the linearised impulse responses with their nonlinear
counterparts (a fixed--Jacobian Newton solve) confines the concern to one
variable. Aggregate output is essentially linear: the ratio of the nonlinear to
the linear peak response, $y^{\mathrm{NL}}/y^{\mathrm{L}}$, stays in $[0.97,1.00]$
across shock sizes, a small nonlinearity comparable to what AMRS report. The
nonlinearity is concentrated in the adoption share $\DG$, whose peak ratio rises
monotonically from $1.13$ at a one-eighth shock to $1.77$ at a three-quarter
shock (Figure~\ref{fig:nonlin}, left). Freezing the margin does not reveal a more
nonlinear economy underneath: under the common--steady-state counterfactual
(Section~\ref{sec:counterfactual}), with $\DG$ held fixed along the transition,
output is itself only mildly nonlinear ($y^{\mathrm{NL}}/y^{\mathrm{L}}\approx
1.15$ at the largest shock), and freezing the margin widens the wedge between the nonlinear and linear paths
rather than narrowing it. The discrete choice is where curvature concentrates, but leaving it free keeps aggregate output within
three percent of its linear approximation.\footnote{At the baseline shock size
($\varepsilon=1.0$) the nonlinear solution does not converge, so the sweep is
reported to $\varepsilon=0.75$. The $\DG$
ratio is monotone and convex over $[0,0.75]$, so nothing depends on the largest
point.}


The convexity is the intrinsic curvature of a logit share evaluated near a low
level. For a binary choice $P=[1+\exp(-\Delta/\sigma_\varepsilon)]^{-1}$ in the
value gap $\Delta$, the relative curvature is $(1-2P)/\sigma_\varepsilon$, so
$D^{G,\mathrm{NL}}/D^{G,\mathrm{L}}-1\propto 1/\sigma_\varepsilon$. Re-solving on
a grid of $\sigma_\varepsilon$ (holding $\bar D^{G}=0.05$ by recalibrating
$\psi_g$) confirms this: the ratio falls monotonically from $1.80$ at
$\sigma_\varepsilon=0.02$ to $1.07$ at $\sigma_\varepsilon=0.40$
(Figure~\ref{fig:nonlin}, right), and it tends to one as the shock vanishes, so
the linearisation is exact to first order. Two
implications follow. Aggregate and consumption-based welfare statistics are
reliable under linearisation. The \emph{level} of induced greening at the
largest shocks is understated by the linear solution, and does not converge
under a full-size shock, so adoption-channel magnitudes are read at a reduced
shock size. The taste scale $\sigma_\varepsilon$ that governs this curvature is
the same parameter Section~\ref{sec:taste_id} and Appendix~\ref{app:taste}
identify: an adoption-elasticity moment pins both at once.

\begin{figure}[H]\centering
\figtitle{Nonlinear against linear responses}
\includegraphics[width=\textwidth]{output/fig9_nonlinearity.png}
\label{fig:nonlin}
\fignotes{Left: peak nonlinear/linear ratio against shock size, for output
(adoption on and off) and $D^{G}$. Right: the same ratio against the logit taste
scale $\sigma_\varepsilon$ at a half-size shock, $\psi_g$ recalibrated to
$\bar D^{G}=5\%$.}
\end{figure}


\subsection{First-order welfare and the nonlinear check}
\label{sec:nlcev}

The consumption-equivalent variations reported throughout are first-order in the
shock: flow utility along the linear impulse response is a first-order deviation
of aggregate felicity, so the CEV is the leading term of the welfare expansion
around a distorted steady state along a deterministic path. The setting is not
the stochastic one in which first-order welfare evaluation is known to give
spurious rankings \citep{KimKim2003,BenignoWoodford2005}: the first-order term is
nonzero and dominant here. But the shock is large and the adoption margin is
convex, so the size of the neglected second-order terms is an empirical question,
which we settle by recomputing each scenario's welfare on the nonlinear
perfect-foresight transition (Table~\ref{tab:nl_cev_import}).

\input{output/tab_nl_cev_import.tex}

Three things follow. First, the linear CEV \emph{overstates} the welfare loss, by
$1$--$4\%$ at a quarter-size shock and up to $7$--$26\%$ at three-quarters, in the
direction the convexity of the adoption margin predicts; the paper's levels are
upper bounds on the losses, and would be overstated by an extrapolated
$10$--$15\%$ at the full shock size, which does not converge nonlinearly. Second,
the policy ordering, flat $>$ Slutsky $>$ cap $>$ ETS $>$ laissez-faire, holds
at every converged size, linear and nonlinear alike, with the first-order gaps
essentially unchanged; the welfare rankings stand. Third, the one difference in the
other direction is the impatient type under the flat transfer, whose
nonlinear/linear ratio rises to $1.89$: the linear solution understates the gain
of high-MPC households when the transfer relaxes a binding liquidity constraint,
which reinforces rather than weakens the distributional result of
Section~\ref{sec:flat}. The single claim the first-order calculation cannot
support is the sign of the crisis-only greening dividend of
Section~\ref{sec:routeA}, which is why we report it there as zero to first order.
